\section{Cơ sở nền tảng và Các hệ thống liên quan}
\label{sec:background}

\noindent \textbf{Text-to-SQL trên lược đồ thực tế và Vòng lặp Tự sửa lỗi.} Bộ chuẩn BIRD~\cite{li2023bird} chứng minh khoảng cách lớn giữa truy vấn mẫu đơn giản và bài toán doanh nghiệp đòi hỏi kết nối nhiều bảng cùng tri thức nghiệp vụ. Các tác tử theo mẫu ReAct~\cite{yao2023react} và Reflexion~\cite{shinn2023reflexion} cải thiện đáng kể tỷ lệ thành công nhờ phản hồi lỗi từ hệ quản trị cơ sở dữ liệu. Tuy nhiên, các nghiên cứu này thường giả định môi trường tin cậy; nếu cho phép sửa lỗi vô điều kiện khi bị từ chối quyền, kẻ tấn công có thể dò đoán cấu trúc hoặc dữ liệu nhạy cảm qua thông báo lỗi. SQLGuard khắc phục điểm yếu này bằng cách chỉ kích hoạt \textit{Bounded Reflection} cho lỗi kỹ thuật (\texttt{repairable}) và ngắt ngay lập tức với vi phạm chính sách (\texttt{terminal}).

\noindent \textbf{Prompt Injection và Giới hạn của các lớp bảo vệ hiện tại.} Greshake và cộng sự~\cite{greshake2023indirect} cùng danh mục rủi ro OWASP LLM~\cite{owasp2025llm} chỉ ra rằng \textit{Indirect Prompt Injection} từ dữ liệu bên ngoài có thể chiếm quyền điều khiển công cụ khi LLM đọc trực tiếp văn bản chưa kiểm dịch. Trong thực tế, các hệ thống phân tích dữ liệu thường vướng một trong ba hạn chế: (1) tin cậy \textit{prompt guardrails} dễ bị qua mặt; (2) áp dụng RBAC mức bảng truyền thống trong DB khiến câu hỏi hỗn hợp bị từ chối toàn phần hoặc buộc phải mở dư cột nhạy cảm; hoặc (3) thiếu giao thức công cụ tách biệt ngữ cảnh vai trò. Vì vậy, SQLGuard kết hợp phân tích cú pháp tĩnh AST qua SQLGlot~\cite{sqlglot} để kiểm soát chi tiết từng biểu thức/cột và chuẩn hóa giao diện công cụ qua Model Context Protocol (MCP)~\cite{mcp2024} với ngữ cảnh vai trò gắn cứng phía máy chủ.
